\section{当前边界与研究路线：从 ``能跑'' 到 ``可靠协商''}

\subsection{为什么自然语言通信是机会窗口}
Brown 指出，多智能体研究历史上长期卡在 ``如何让体之间形成通信协议''。LLM 时代这一步骤被显著简化，Agent 可直接用自然语言协商，这是前所未有的红利。

\begin{importantbox}{时代性变化}
过去需要专门学习 emergent communication，现在默认就有高带宽语言通道。研究瓶颈从 ``会不会说'' 转向 ``说了能否可靠达成一致''。
\end{importantbox}

\subsection{行业复盘：为什么很多团队说 ``先别急着多 Agent''}
课程引用了 Cognition 与 Anthropic 的工程复盘：多 Agent 在部分场景有效，但系统常见特征是 ``有效但脆''。典型问题包括：
\begin{itemize}
    \item 长上下文协商中目标漂移；
    \item 多轮互审导致延迟和成本失控；
    \item 冲突协调失败时，整体可靠性不如单 Agent。
\end{itemize}

\begin{knowledgebox}{可落地的可靠性指标}
\begin{itemize}
    \item 协议一致率：多 Agent 对任务目标和约束解释是否一致；
    \item 冲突收敛率：分歧是否在有限轮内收敛；
    \item 代价回报比：相对单 Agent 的质量提升是否覆盖额外 token 与时延；
    \item 失败可恢复性：单个子 Agent 异常是否会拖垮全局。
\end{itemize}
\end{knowledgebox}

\subsection{从课程结论到工程决策}
前几节分别讨论了自然语言通信的机会和现有 scaffold 的脆弱性，现在需要把它们压缩成可执行的选择顺序。这里不再问“要不要上多 Agent”，而是先辨认任务的博弈结构与目标群体，再判断并行采样、专家路由或长程协商中的哪一种机制值得承担额外成本。
把本讲压缩成一句决策规则是：\textbf{先判定任务结构，再选优化目标，再选算法家族}。
如果任务更接近二人零和，优先 exploitability 与稳健收敛；
如果任务是多人协作，优先群体分布建模与跨群体验证。

\begin{warningbox}{部署前的最低防线}
不要把 ``demo 上可运行'' 误判为 ``生产可依赖''。多 Agent 系统必须在对抗输入、跨分布人群、长任务链条下做稳定性回归。
\end{warningbox}

\subsection{本章小结}
LLM 让多 Agent 研究从 ``通信可行性'' 跨到 ``协商可靠性'' 阶段。接下来几年的主战场是评测协议、鲁棒协同机制与成本可控的系统化落地。

